% 表1：模型与两阶段设计。全部共用M=64、噪声标准差0.05、B=500及α=0.05。
\begingroup
\fontsize{8.5}{11}\selectfont
\setlength{\tabcolsep}{3pt}
\renewcommand{\arraystretch}{1.15}
\begin{longtable}{p{18mm}p{31mm}p{31mm}p{32mm}p{32mm}}
\toprule
模型 & 潜在背景 & 信号结构 & 探索形状 & 确认形状 \\
\midrule
\endfirsthead
\toprule
模型 & 潜在背景 & 信号结构 & 探索形状 & 确认形状 \\
\midrule
\endhead
M1 & 五基傅里叶；权重随基序平方倒数衰减 & 与背景同形的跨通道协方差 & (200,20)、(400,20)、(800,20)、(800,50) & (800,20) \\
M2 & 规范化布朗运动；63个增量 & 与背景同形的跨通道协方差 & (200,20)、(400,20)、(800,20)、(800,50) & (800,20) \\
coherent & 20维平坦谱 & 秩一相干信号；中心3/8 & (200,20)、(400,20)、(800,20)、(800,50) & (400,20)、(800,20)、(800,50) \\
local & 20维平坦谱；共享flat H0 & 原Fejér局部信号 & (800,20) & 无 \\
fast & 20维正余弦成对快速衰减谱 & 原冻结fast信号 & (800,20) & 无 \\
\bottomrule
\end{longtable}
\endgroup
\begingroup
\fontsize{8.5}{11}\selectfont
\setlength{\tabcolsep}{3pt}
\renewcommand{\arraystretch}{1.15}
\begin{longtable}{p{23mm}p{30mm}p{25mm}p{32mm}p{32mm}}
\toprule
阶段 & H0/H1单元 & 数据集 & 方法/配对 & H0/H1重复 \\
\midrule
\endfirsthead
\toprule
阶段 & H0/H1单元 & 数据集 & 方法/配对 & H0/H1重复 \\
\midrule
\endhead
探索 & 13/56 & 9300 & 276/168 & 500/50 \\
独立确认 & 5/15 & 9500 & 80/45 & 1000/300 \\
\bottomrule
\end{longtable}
\endgroup
\noindent{\footnotesize 探索$\tau$=0/0.15/0.30/0.45/0.60，确认$\tau$=0/0.15/0.30/0.60；$\tau$=0为H0，其他为H1；local/fast仅在anchor探索。flat H0只计一次。}
